\section{Finite critical histories and physical readouts}
\label{sec:bulk}

The entry construction provides a small interaction at a large scale, but
its cutoff is useful only on a finite range of larger volumes.  We therefore
work with a family of finite histories.  This section proves that these
histories stay close to the marginal Gaussian reference, and obtains two
physical observables that compare histories with different entry scales.
Section~\ref{sec:trajectory} will use those comparisons to determine the
gradient coefficient with a summable error.

Throughout this section $a=\alpha^2=8\pi$, and $L$ is the fixed block
factor chosen after the geometric and analytic norm parameters.  An
\emph{absolute index} $k$ refers to the physical side $L^k$, rather than
to the number of steps since entry.  This distinction matters when two
histories have different entry scales.

\subsection{The family of histories and the local maps}

Fix the parameters of Proposition~\ref{prop:entry}, including $R$, $P=R^5$
and $q$, independently of the history length.  For a large integer $J$,
choose a dyadic observation side $s=s_J$ such that
\begin{equation}\label{eq:bulk-entry-scale}
 J^3\le s_J<C_LJ^3,\qquad qRs_J=L^{i(J)}.
\end{equation}
Such a choice is possible because $q,R,L$ are powers of two: among any
fixed number of successive dyadic choices of $s$, one has the required
congruence of exponents.  In particular $i(J)=O_L(\log J)$.  Starting
from the unpierced plane activity of Proposition~\ref{prop:entry}, run
the local maps through absolute index $8J$, stopping provisionally if
the activity leaves the small ball of the map theorem.  We call this
the $J$-history.  A superscript $(J)$ distinguishes its coordinates
when necessary.
Figure~\ref{fig:histories} displays the overlap of two such histories
and the intervals that will be retained when they are joined in
Section~\ref{sec:trajectory}.

\begin{figure}[tb]
\centering
\begin{tikzpicture}[x=0.65cm,y=0.7cm,>=Latex,font=\small]
 \draw[->] (0,0)--(16.8,0) node[right] {scale index};
 \foreach \x/\lab in {2/J,4/2J,8/4J,16/8J}
  {\draw (\x,0.08)--(\x,-0.08) node[below] {$\lab$};}
 \draw[thin] (0.55,3)--(16,3);
 \draw[very thick,blue!55!black] (2,3)--(4,3);
 \draw (0.55,2.87)--(0.55,3.13);
 \draw (16,2.87)--(16,3.13);
 \node[above] at (0.8,3.15) {$i(J)=O(\log J)$};
 \node[above] at (8.7,3.08) {history targeted at $J$};
 \draw[thin,->] (0.82,1)--(16.3,1);
 \draw[very thick,blue!55!black] (4,1)--(8,1);
 \draw (0.82,0.87)--(0.82,1.13);
 \node[above left] at (16,1.08) {history targeted at $2J$ (to $16J$)};
 \draw[dashed] (4,0.2)--(4,3.3);
 \draw[->] (4.3,2.78)--(4.3,1.23);
 \node[right] at (4.5,2.05) {compare physical observables};
 \node[below,align=center] at (8,-0.8)
  {Thick segments are retained in the spliced sequence.};
\end{tikzpicture}
\caption{Finite histories overlap well beyond the intervals retained in
the spliced sequence. Coordinates are compared at the same absolute
scale using observables of the original height law. The picture is
schematic near the entry indices; no convergence rate at entry is
assumed.}
\label{fig:histories}
\end{figure}


We use the singleton and remainder coordinates from
Proposition~\ref{inp:balanced-map}.  The unbalanced coordinates are
$(t,z,\mathcal R)$: the first two multiply the gradient singleton
$e_B^0/2$ and the fundamental singleton $c_B^\alpha$, respectively.
The balanced gradient and fundamental coordinates are $(T,Z)$.
Parametrizing the exact Gaussian curves gives the coordinate $u$, and
subtracting the corresponding Gaussian remainder gives $B=B_e+B_o$,
where the subscripts denote the even and odd parts under translation
by $\pi/\alpha$.  In these coordinates an exact Gaussian trajectory
is $(u,0,0,0)$ with $u$ unchanged by the map.  The curve depends on its
entry index; the reference covariance kernels at a specified absolute
index do not.

\begin{lemma}[Uniform maps and the absolute coefficient estimates]
\label{lem:bulk-parity}
The coordinate changes and parity-resolved maps of
Proposition~\ref{inp:balanced-map} have a common analytic neighborhood
and common bounds for all sufficiently large entry indices $i(J)$.
At absolute index $k$, their two scalar quadratic coefficients obey
\begin{equation}\label{eq:bulk-coefficient-rate}
 H_k=H+O((1+k)^{-2}),\qquad
 b_k=b_*+O((1+k)^{-2}),\qquad
 H>0,\quad b_*=2\log L.
\end{equation}
The coefficients $H_k,b_k$ are independent of the entry index.
The constants in the parity remainder bounds are uniform in the
$J$-history.
\end{lemma}

\begin{proof}
The balanced projections are the convergent future linear projection
series in \cite[Equations (8.17) and
(8.18)]{B}.  Their norms and inverse norms are uniform
at every sufficiently large scale.  The covariance used here is the
nearest-neighbor polynomial covariance in the fixed absolute-index
convention.  Consequently its linear maps and future projection
series, and hence its balanced quadratic coefficients, are functions
of the absolute index alone.

We first check that parametrizing the Gaussian curves does not
introduce a dependence of the analytic radius on the entry scale.
At side $L^i$, initialize the exact partition algebra with the
exponential of $d\sum_B e_B^0/2$.  A fixed small complex disk in $d$
is permitted by the energy reserve in the entry regulator.  The
normalized block energy has the same bound at every entry scale.
For a smooth terminal source, Gaussian completion involves the
inverse of $I-dC^{1/2}AC^{1/2}$, where $C$ is the reference tail
covariance and $A$ is the unit-edge Laplacian.  Energy contraction
bounds this inverse uniformly on a fixed smaller disk.  The
gradient diagnostic and transverse contraction in
\cite[Lemma 10.3]{B} therefore bound all these
Gaussian trajectories and their analytic derivatives uniformly.
Smaller diagnostic tori cause no problem in this argument: the
Gaussian initialization has no entry sewing operations.  This is
also the uniformity argument used at the end of
\cite[Section 5]{X}.  The cosine decomposition
does not change it, since the pure Gaussian trajectory has zero
charged coordinates.  Its balanced gradient coordinate is
$T_k^g(d)=d+O(d^2)$, with uniform first derivative, and therefore
has an inverse on a common disk.  The transverse derivative of the
Gaussian remainder at $d=0$ is zero.  The parity estimates of
\cite[Lemma 10.6]{B} thus hold uniformly after this
change of coordinates.  The change leaves the $Z^2$ and $uZ$
coefficients unchanged.

The limits and signs in \eqref{eq:bulk-coefficient-rate}, with
$b_*=\alpha^2\log L/(4\pi)=2\log L$, follow from
\cite[Proposition 8.5 and Equation
(8.21)]{B}.  We need a rate in this particular
calculation.  We give the quantitative version of its finite-depth
approximation, since qualitative convergence of these coefficients
would not suffice below.

Truncate each future projection series at depth $r$.  Uniform
contraction of the linear remainder map makes the discarded part
at most $C\vartheta^r$ for some $\vartheta<1$.  Shell eigenvalue
errors and the kernel errors in fixed-order scaled differences are
power-small in $L^k$, by \cite[Lemma 8.1 and Equation
(8.9)]{B}.  We claim that replacing all kernels in
the retained depth-$r$ calculation by their continuum kernels has
error bounded by
\begin{equation}\label{eq:bulk-finite-depth-error}
 \exp\{C_L(1+r)^p\}L^{-ck}
\end{equation}
for fixed $p<\infty$ and $c>0$.

Indeed, a quadratic image starts with two singleton inputs and has
bounded support size.  Thereafter only linear maps occur.  A linear
continuation takes block closures, extracts a singleton, or
subtracts that extraction.  Support size consequently remains
bounded, and the number of shapes, routed blocks and extraction
choices through depth $r$ is at most
$C_L^{O((1+r)^2)}$.  Write singleton energies throughout in scaled
differences with their volume-normalized cell sums.  The original
two inputs have bounded polynomial degree.  Each subsequent Taylor
projection takes at most two derivatives on a constant or affine
background, with output a quadratic polynomial, a constant, or a
single phase harmonic.  With independent variables for the successive
Gaussian shells, the resulting Gaussian polynomial-character
expectations have at most $C_L(1+r)$ polynomial or derivative slots.
Their pairing expansions have at most factorial growth in this
number of slots.  Scale ratios and distances are bounded by
$C_L^{1+r}$.  The real character exponential has modulus at most
one; its change under kernel replacement can be estimated along
the line segment between two positive covariance matrices.  These
facts bound the expectation, and its derivative with respect to
any covariance entry, by the exponential factor in
\eqref{eq:bulk-finite-depth-error}.

Only values and fixed-order scaled differences of kernels enter
these slots.  Later Taylor projections differentiate the remaining
background field, not previously integrated kernels.  Converting
between the finitely many scales introduces at most
$C_L^{O((1+r)^2)}$.  The volume-normalized lattice sums satisfy the
same bound: couple each lattice sample to its mesh cube and use
one additional derivative of the smooth continuum shell kernel.
The block-shape sums are unchanged by this coupling.  This proves
\eqref{eq:bulk-finite-depth-error}; in particular it is a bound on
the explicit singleton calculation, not an assertion of uniform
Riemann approximation for arbitrary activities.

Choose $r=C_0\log(2+k)$.  A sufficiently large $C_0$ makes the
projection tail $O((1+k)^{-3})$, while
$\exp\{C_L(1+r)^p\}L^{-ck}=O((1+k)^{-3})$ for large $k$.
Weakening the bound gives \eqref{eq:bulk-coefficient-rate}.
\end{proof}

\subsection{Periodic physical diagnostics}

We now identify observables which test the two nearly marginal
coordinates.  Fix a sufficiently large dyadic $D_0$ and use tori of
side
\begin{equation}\label{eq:bulk-stopping-tori}
 N_k=D_0L^k.
\end{equation}
The constant $D_0$ exceeds the causal clearance of the local maps.
For each admissible stop, the torus has its own exact recursion from
the same entry scale as the plane history.  Its singleton coefficients
agree with the plane ones.  Its remainder need not agree with the
plane remainder and will be estimated separately.

Only stops with $N_k\ge sP$ are used.  Before this condition holds
there are at most $O(1+\log R)$ plane steps.  Their bounded map
amplification is $R^{O(1)}$, which is absorbed by the exponentially
small entry norm when $R$ is increased.  All stopping tori are
compatible with the observation and tile grids.

Let $f$ be the first real horizontal cosine on the unit torus, and
put
\begin{equation}\label{eq:bulk-test-definitions}
 g=(f,(-\Delta_{\mathbb T^2})^{-1}f)>0,\qquad
 X_N=(y,f_N),\qquad \chi_N=e^{i\bar y}.
\end{equation}
Here $f_N$ is sampled at observation-cell centers and includes the
cell-area weights, and $\bar y$ is the uniform mean of the
observations.  Heights are taken modulo a common translation by
$2\pi$, so $e^{i\bar h}$ is well defined.  The expectation in the
next lemma is for the periodic Bessel height law before inserting
the cutoff, together with its independent observation noises.

\begin{lemma}[The two physical limits]
\label{lem:bulk-physical-diagnostics}
As $N\to\infty$ with $s/N\to0$, locally uniformly for complex $w$,
\begin{align}
 \E e^{wX_N}&\longrightarrow e^{agw^2/2},
 \label{eq:bulk-physical-gradient}\\
 \E(\chi_Ne^{wX_N})-\E\chi_N\,e^{agw^2/2}
 &\longrightarrow0.
 \label{eq:bulk-physical-phase}
\end{align}
All fixed source derivatives converge as well.  For $s/N\le P^{-1}$,
the discrepancies can be made uniformly small by first increasing
$R$, hence $P$, and then increasing $s$.  The same statements hold
with the cutoff on every stopping torus of a $J$-history, uniformly
up to index $8J$.
\end{lemma}

\begin{proof}
We first transfer the height limit in
\cite[Theorem 2.4]{X} from free rectangles to
the periodic graph.  For a smooth periodic function $v$, test the
height by $(\partial h,\partial v_N)$.  Cut a fixed macroscopic
grid of seams to obtain free rectangles.  The retained tests depend
only on their internal gradients.  The centered comparison on the
independent Gaussian chains therefore bounds their variances and
real transforms by those on the cut graph.  There are $O(N)$ seam
edges, with test coefficients $O(N^{-1})$.  Their squared flow
cost is $O(N^{-1})$, so the omitted test tends to zero in reference
cost by \cite[Lemma 2.3]{X}.  The free
rectangle limits give the upper bound
$a\int|\nabla v|^2$ for each subsequential covariance limit.

For a function $v$ supported in a coordinate rectangle, choose a
pinning band outside its support and ground one representative
height at a point of the band.  Adding the hard pin decreases the
variance by the finite comparisons.  The mixed pinned limit in
\cite[Theorem 3.1]{X} gives the lower bound
$a\int|\nabla v|^2$, since $v$ is supported strictly inside the
pinning band.  Hence equality holds for such local functions.
The upper-bound defect of any subsequential covariance limit is
a positive semidefinite form.  A vector on which its quadratic
form vanishes belongs to its null space, by Cauchy--Schwarz for
positive semidefinite forms.  A smooth partition of unity writes
every periodic $v$ as a finite sum of these local functions.
Equality thus holds for every $v$, including their mixed
covariances.  The comparison argument is valid on the periodic
Gaussian-chain approximation at a fixed graph before passing to
the Bessel limit; no bare inverse-energy bound is transferred.

The corresponding Laplace limits follow by the same two comparisons.
For the lower bound, the centered variance comparison used in the
proof of \cite[Theorem 2.4]{X} supplies the
variance-based lower bound for the transform.  For the upper bound,
cut to the fixed grid of free rectangles and use their Laplace
limits.  H\"older's inequality separates the omitted seam term;
its reference cost tends to zero, and the H\"older exponent on
the retained term can then be sent to one.  Uniform exponential
bounds from \cite[Lemma 2.3]{X} give
convergence of every fixed moment and local uniform convergence
of the complex transforms.  Smooth mean-free density tests follow
by solving the periodic Poisson equation and approximating in
reference energy.  For the first cosine this also follows directly
from its discrete Laplacian eigenvalue.

Cell sampling changes the normalized reference standard deviation
of this test by $O(s/N)$, by the periodic Poincar\'e inequality
or the cell flow estimate.  Its observation-noise standard
deviation is also $O(s/N)$.  This proves
\eqref{eq:bulk-physical-gradient}, including its differentiated
version.

The phase requires more than the marginal Gaussian limit.  We
adapt the replica calculation of
\cite[Lemma 10.1]{B}.  Remove the
independent observation noise temporarily, and write $X_N^{(h)}$
for the height part of the cosine test.  We first work at one
fixed Gaussian-chain approximation of the size-$N$ graph.  The
pin selects a height representative and its law is invariant
under $h\mapsto-h$.  Let $h_1,h_2$ be independent copies on
the entire subdivided graph, including the chain variables, and put
\[
 S=h_1+h_2,\qquad \Delta=h_1-h_2.
\]
Their common coordinatewise parity class is
$\mathcal P=S\bmod4\pi\Z=\Delta\bmod4\pi\Z$.  If $Q$ denotes
the quadratic precision form, the change of variables gives
$Q(h_1)+Q(h_2)=(Q(S)+Q(\Delta))/2$.  Conditional on
$\mathcal P$, the fields $S,\Delta$ are consequently independent
and identically distributed.  Their conditional laws are centered:
each parity coset is preserved by sign reversal.  This is the
replica identity of \cite[Equation (2.12), in the
proof of Theorem 2.5]{X}.

For this fixed approximation write $X(h)$ for the cosine test and
$\mu(h)$ for the mean over the original sites, both extended as
linear forms on the subdivided field.  Thus $X(h)$ approximates
$X_N^{(h)}$ and $\mu(h)$ approximates $\bar h$.  Define the scalar
sum and difference tests and their conditional variance by
\begin{align*}
 A&=X(h_1)+X(h_2)=X(S),\\
 D&=X(h_1)-X(h_2)=X(\Delta),\\
 V&=\E(A^2\mid\mathcal P)=\E(D^2\mid\mathcal P).
\end{align*}
Put $\sigma^2=\E X(h)^2$ and
$\kappa_4=\E X(h)^4-3\sigma^4$.  Conditional independence gives
$\E V^2=\E(A^2D^2)$.  Since
$A^2D^2=(X(h_1)^2-X(h_2)^2)^2$, independence of the original
replicas yields $\E V^2=2\E X(h)^4-2\sigma^4$ and therefore
\begin{equation}\label{eq:bulk-replica-variance}
 \E V=2\sigma^2,\qquad \operatorname{Var}V=2\kappa_4.
\end{equation}

For real $t$ define the transform and the difference-character
\[
 F(t)=\E e^{i\mu(h)+itX(h)},\qquad
 \Psi_t= e^{i\mu(\Delta)+itX(\Delta)}.
\]
The character $\Psi_t$ depends only on the difference field and
has modulus one.  Conditional independence gives
$\E(A^2\Psi_t)=\E(V\Psi_t)$.  Inversion symmetry of the second
replica gives $\E\Psi_t=F(t)^2$.  The same sign reversal and
differentiation of $F$ give, with primes denoting $t$ derivatives,
\[
 \E(A^2\Psi_t)
 =(-F'')F+F(-F'')+2(-iF')(iF')
 =-2\bigl(FF''-(F')^2\bigr).
\]
We have thus proved the exact finite-approximation identity
\begin{equation}\label{eq:bulk-weighted-replica}
 -2\bigl(FF''-(F')^2\bigr)
 =2\sigma^2F^2+
   \E\bigl[(V-2\sigma^2)\Psi_t\bigr],
 \qquad t\in\R.
\end{equation}
Cauchy--Schwarz and \eqref{eq:bulk-replica-variance} bound its
last term by $\sqrt{2\kappa_4}$.

Take the chain limit first, at the fixed original graph.  The
fixed-graph moment convergence passes the second and fourth
moments, the transforms and their first two derivatives through
this limit.  Consequently, for
$F_N(t)=\E e^{i\bar h+itX_N^{(h)}}$ in the original Bessel law,
the absolute difference between the two sides of
\eqref{eq:bulk-weighted-replica} without its last term is at most
$\sqrt{2\kappa_{4,N}}$, where
$\kappa_{4,N}=\E(X_N^{(h)})^4-3(\E(X_N^{(h)})^2)^2$.
The Gaussian moment convergence proved above gives
$\kappa_{4,N}\to0$ and
$\operatorname{Var}(X_N^{(h)})\to ag$.
Uniform exponential bounds on compact complex sets permit
locally uniform subsequential limits of $F_N$ and all their
fixed derivatives.  For any such limit, denoted again by $F$,
\eqref{eq:bulk-weighted-replica} therefore implies on the real
axis, and then everywhere by analytic continuation,
\begin{equation}\label{eq:bulk-replica-ode}
 FF''-(F')^2=-agF^2.
\end{equation}
Half-wavelength translation preserves $e^{i\bar h}$ and changes
the sign of the cosine test.  Repinning after translation only
subtracts a height in $2\pi\Z$, so it leaves this phase unchanged.
Hence $F'(0)=0$.  If $F$ is not identically zero, solving the equation
where $F\ne0$ and then continuing analytically shows that $F$ is
an exponential quadratic without zeros.  Thus
$F(t)=F(0)e^{-agt^2/2}$.  The identically zero limit has the same
factorization.  Restoring the independent observation noise proves
\eqref{eq:bulk-physical-phase}.

These arguments also give the asserted order of choices.  Compare
the cell test to the ordinary sampled cosine with an error bounded
by $C/P$ when $s/N\le P^{-1}$.  First make this bound small, and
then use convergence of the ordinary cosine tests for all
$N\ge sP$.  Finally the cutoff error
\eqref{eq:cutoff-error}, including fixed exponential tests and
source derivatives, is bounded throughout a $J$-history by
\[
 C\exp\{C_LJ-cs_J\}=O(e^{-c'J^3}).
\]
Consequently it changes none of the preceding conclusions.
\end{proof}

We record the map interpretation of these tests.  A source $wX_N$
translates the remaining field by a smooth function $F$ satisfying
$B_s\alpha F$ equal to the reference observation covariance applied
to the source.  It can be chosen as a microscopic cosine: translation
symmetry makes the observation-cell function a cosine, and cell
averaging of the corresponding microscopic cosine has a multiplier
bounded away from zero.  Gaussian completion gives
\begin{equation}\label{eq:bulk-F-energy}
 (F,AF)=aw^2g+O_w((s/N)^c+N^{-c})
\end{equation}
for a fixed $c>0$, with uniform scaled derivatives in the point
norm.  One can obtain this error by periodic inverse-energy
approximation of the sampled density and the Taylor expansion of
the first discrete eigenvalue.

At a stop with $D_0^2$ blocks, the analytic terminal expansion of
\cite[Equations (10.3) and
(10.4)]{B} gives
\begin{align}
 \log\E e^{wX_N}-\log\E_a e^{wX_N}
 &=\tfrac12t_k(F,AF)
   +O(\norm{\mathcal R_k^{\rm tor}}+\nu_k^2),
 \label{eq:bulk-map-gradient}\\
 e^{-\log\E_a e^{wX_N}}\E(\chi_Ne^{wX_N})-\E\chi_N
 &=z_k\mathcal M_k(w)
   +O(\norm{\mathcal R_k^{\rm tor}}+\nu_k^2).
 \label{eq:bulk-map-phase}
\end{align}
Here $\nu_k=|t_k|+|z_k|+\norm{\mathcal R_k^{\rm tor}}$,
and $\E_a$ denotes the reference Gaussian observation law.  For
one fixed sufficiently small nonzero real $w$, both displayed
linear multipliers are bounded away from zero.  Explicitly the
second multiplier is
\begin{equation}\label{eq:bulk-phase-multiplier}
 \mathcal M_k(w)=\frac{D_0^2}{2}
 e^{\tau_k/2-\alpha^2D_{{\rm tail},k}/2}
 \operatorname{ave}_x(e^{-i\alpha F(x)}-1).
\end{equation}
The tail diagonal $D_{{\rm tail},k}$ is bounded, and the last
average tends to a strictly negative number for the fixed $w$.
The number $\tau_k$ is the variance, in observation units, of
the integrated Gaussian means.  It is bounded because the
zero-mode continuation at side $u$ contributes $O(u^2/N_k^2)$.
The compensating factor $e^{\tau_k/2}$ is necessary: insertion
of the retained uniform mean gives the physical mean-phase
insertion times $e^{-\tau_k/2}$.  This prescription is exact for
the cutoff interaction, which is periodic in the common mean.
Before tilt and subtraction, the same formula with spatial
factor one is a real positive multiplier bounded above and below.

\subsection{No exit and the nonexponential odd sector}

\begin{lemma}[Smallness on the finite horizons]
\label{lem:bulk-window-vanishing}
For every sufficiently small prescribed tube radius, the entry
parameters can be chosen so that all $J$-histories, for sufficiently
large $J$, remain in that tube through index $8J$.  The stopped
torus remainders remain in a fixed comparable tube.  Moreover,
for every fixed $0<c<8$,
\begin{equation}\label{eq:bulk-window-vanishing}
 \sup_{cJ\le k\le8J}
 \bigl(|u_k^{(J)}|+|Z_k^{(J)}|
       +\norm{B_{e,k}^{(J)}}+\norm{B_{o,k}^{(J)}}\bigr)
 \longrightarrow0.
\end{equation}
The analogous stopped torus remainder quantities tend to zero.
\end{lemma}

\begin{proof}
We use the two diagnostics exactly at the finite stops at which
they have just been justified.  Before a first candidate exit,
contraction of the plane remainder, and separately contraction of
the remainder on each stopping torus, give
\[
 \norm{\mathcal R_k},\ \norm{\mathcal R_k^{\rm tor}}
 \le C\vartheta^{k-i(J)}\varepsilon+C\rho^2.
\]
Here $\rho$ bounds the coordinates before the candidate exit,
$\varepsilon$ is the entry norm, and $\vartheta<1$.  Boundedness
of one map keeps the candidate state in a fixed multiple of that
tube.  Equations \eqref{eq:bulk-map-gradient} and
\eqref{eq:bulk-map-phase} and the nonzero multipliers bound its
singleton coordinates by
$C(\delta_{\rm diag}+\varepsilon+\rho^2)$, where
$\delta_{\rm diag}$ bounds the physical discrepancies.
Choose $\rho$ so that the quadratic term is smaller than a
fixed small multiple of $\rho$; next choose $R$ large and the
entry norm and diagnostic discrepancies small compared with
$\rho$.  Lemma~\ref{lem:bulk-physical-diagnostics} and
Proposition~\ref{prop:entry} permit these choices.  For large
$J$ the candidate cannot be an exit.  The finitely many initial
plane-only steps were controlled after
\eqref{eq:bulk-stopping-tori}.  The coordinate changes are
uniformly bounded by Lemma~\ref{lem:bulk-parity}, so this argument
applies to either coordinate system.

For the convergence assertion, take the supremum limiting plane
size over sequences with $J\to\infty$ and
$k-i(J)\to\infty$, $k\le8J$.  Call it $q_0$.  It lies in the
small tube, and all fixed numbers of preceding steps satisfy
the same limiting bound.  Iterating contraction over $r$ steps
and then sending $r\to\infty$ bounds the limiting plane
remainder by $Cq_0^2$.  The torus remainder has its own
contraction with the same plane singleton forcing, so the
same bound holds at its separate stops.  The test errors tend
to zero because $s/N_k\to0$ on these sequences.  The diagnostics
therefore give $q_0\le C'q_0^2$.  Taking the original tube
smaller than $1/C'$ implies $q_0=0$.  Every window in
\eqref{eq:bulk-window-vanishing} satisfies
$k-i(J)\to\infty$, which proves the stated uniform result.
\end{proof}

The preceding lemma gives convergence to the Gaussian fixed point.
To determine the approach rate we must also rule out exponential
decay in the odd direction.  The relevant lower bound is a
statement about the original periodic model at absolute side
$N_k$; it therefore does not depend on the entry history.

\begin{lemma}[Positive odd coordinate and contracted remainders]
\label{lem:bulk-odd}
On every fixed window $cJ\le k\le8J$, $0<c<8$, and for large
$J$, the period convention may be kept unchanged and
\begin{equation}\label{eq:bulk-odd-estimates}
 Z_k^{(J)}>0,\qquad -\log Z_k^{(J)}=o(J),\qquad
 \norm{B_{e,k}^{(J)}}\le C(Z_k^{(J)})^2,
 \qquad
 \norm{B_{o,k}^{(J)}}\le
 C\bigl(|u_k^{(J)}|+(Z_k^{(J)})^2\bigr)Z_k^{(J)}.
\end{equation}
The $o(J)$ and the constants are uniform on the window.
\end{lemma}

\begin{proof}
Let $Y_k=\E e^{i\bar h_{N_k}}$ for the original periodic height
law.  We first show
\begin{equation}\label{eq:bulk-phase-lower}
 Y_k>0,\qquad -\log Y_k=o(k).
\end{equation}
For free squares let $z_n^{\rm face}$ be the positive
characteristic ratio of the averaged four-face profile in
\cite[Equation (2.15)]{X}.  At $a=8\pi$ there is
$c_0>0$, independent of the fixed sufficiently large integer $r$,
such that
\[
 \liminf_{n\to\infty}
 \frac{z_{L^rn}^{\rm face}}{z_n^{\rm face}}\ge c_0.
\]
Replace $c_0$ by a smaller number in $(0,1)$.  For each fixed $r$,
iterate the eventual bound $c_0/2$ on the $r$ residue classes of
the absolute index.  This gives
\[
 \limsup_{k\to\infty}\frac{-\log z_{N_k}^{\rm face}}k
 \le\frac{\log(2/c_0)}r.
\]
Sending $r\to\infty$ proves subexponential decay of this free
characteristic.  The difference of the face profile and the
uniform profile has bounded reference cost.  The phase transfer
inequality \cite[Equation (2.2)]{X}, or its
application in \cite[Lemma 10.7]{B},
changes a logarithmic loss $o(k)$ by at most $O(1+\sqrt{o(k)})$.
Periodic edge restoration increases the positive characteristic
by \cite[Lemma 2.2]{X}.  This proves
\eqref{eq:bulk-phase-lower}.  Observation noise multiplies it
by $\exp\{-O((s/N_k)^2)\}$, and the cutoff changes it by
$O(e^{-cJ^3})$ on the windows under consideration.  Both changes
are negligible compared with \eqref{eq:bulk-phase-lower}.

Fix a slightly larger window starting at a smaller positive
fraction of $J$.  By Lemma~\ref{lem:bulk-window-vanishing} all
coordinates there tend uniformly to zero.  For a fixed small
$\epsilon>0$, consider the cone
\[
 \norm{B_o}\le\epsilon|Z|.
\]
The parity map makes it forward invariant once the window is
sufficiently late.  Indeed inside the cone its scalar equation
gives $|Z_+/Z-1|=o(1)$ uniformly, whereas its odd remainder
equation gives
$\norm{B_{o,+}}\le(\vartheta\epsilon+o(1))|Z|$.
In particular the sign of $Z$ cannot change inside the cone.

While the trajectory lies outside the cone,
$|Z|\le\epsilon^{-1}\norm{B_o}$; the odd remainder equation
then contracts $B_o$ with some factor $\vartheta'<1$.  It gives
the same geometric upper bound on $Z$ until the cone is entered.
This alternative cannot persist for a positive fraction of $J$.
To see this with the physical observable, the stopped mean-phase
expectation has an odd factor under the common-mean integral.
It is bounded by the plane odd singleton and the stopped torus
odd remainder.  The balance relations bound the unbalanced
plane odd variables by $|Z|+\norm{B_o}$.  The torus odd remainder
satisfies a geometric convolution of these plane odd inputs,
plus a geometric transient: every nonlinear odd-sector term
contains an odd factor.  An interval of length proportional to
$J$ outside the cone would therefore make the phase exponentially
small at its last stop, contradicting
\eqref{eq:bulk-phase-lower}.  Taking successively enlarged
windows gives cone entrance strictly before the desired window.

Inside the cone, the adjacent $|Z|$ ratios tend uniformly to one.
For any prescribed small $\delta>0$, the preceding convolution
is thus bounded by $C|Z_k|+Ce^{-cJ}$, since earlier inputs are
at most $(1+\delta)^r|Z_k|$ after $r$ backward steps and
$\vartheta(1+\delta)<1$.  Equation
\eqref{eq:bulk-phase-lower} consequently implies
$|Z_k|\ge e^{-o(J)}$ throughout the desired window.  In
particular all homogeneous contraction transients are negligible
relative to every fixed positive power of $|Z_k|$.

Iterate the even remainder equation with these adjacent ratio
bounds.  Its forcing is $O(Z^2)$, and the transient is negligible,
so $\norm{B_{e,k}}\le C Z_k^2$.  The scalar equation then gives
$|u_{k+1}-u_k|\le CZ_k^2$.  For $r$ steps backward within an
enlarged window,
\begin{equation}\label{eq:bulk-backward-u}
 |Z_{k-r}|\le(1+\delta)^r|Z_k|,\qquad
 |u_{k-r}|\le |u_k|+Cr(1+\delta)^{2r}Z_k^2.
\end{equation}
Using these bounds in the odd contraction gives
\[
 \norm{B_{o,k}}
 \le C\sum_{r\ge0}\vartheta^r
 (|u_k|+C(1+r)(1+\delta)^{2r}Z_k^2)
 (1+\delta)^r|Z_k|
 \le C(|u_k|+Z_k^2)|Z_k|,
\]
where $\delta$ is chosen so that
$\vartheta(1+\delta)^3<1$.  The finitely truncated convolution
has only a negligible transient.  This proves the norm bounds.

Finally the untilted mean-phase expansion, with its positive
linear multiplier described after
\eqref{eq:bulk-phase-multiplier}, is
$M_kZ_k+o(|Z_k|)+O(e^{-cJ})$, where $M_k$ is bounded positively
above and below.  The torus odd remainder has the sharper bound
just obtained by the same contraction argument.  The physical
phase is positive by \eqref{eq:bulk-phase-lower}, and its
cutoff error is negligible.  Thus $Z_k>0$ in the unchanged period
convention.  Since the coordinates also tend to zero, the
lower bound $Z_k\ge e^{-o(J)}$ is exactly the logarithmic
assertion in \eqref{eq:bulk-odd-estimates}.
\end{proof}

\subsection{Readouts common to different entry scales}

The estimates obtained so far do not yet identify a single
trajectory: each $J$-history uses different observation cells and
different cutoffs.  We now compare their coordinates using
observables of the same original periodic model.

Let $\widetilde X_{N_k}$ be the ordinary microscopic sampled cosine,
without observation averaging or noise, and define
\begin{equation}\label{eq:bulk-common-observables}
 U_k=\frac{\operatorname{Var}(\widetilde X_{N_k})}{g},
 \qquad Y_k=\E e^{i\bar h_{N_k}}.
\end{equation}
Both quantities depend only on the absolute index $k$.

\begin{proposition}[Common readouts]
\label{prop:bulk-readouts}
For every $0<c<8$, uniformly on $cJ\le k\le8J$,
\begin{align}
 U_k&=\frac{a}{1-u_k^{(J)}}+O((Z_k^{(J)})^2)+O(e^{-c'J}),
 \label{eq:bulk-U-readout}\\
 Y_k&=M_kZ_k^{(J)}
 +O\bigl((|u_k^{(J)}|+(Z_k^{(J)})^2)Z_k^{(J)}\bigr)
 +O(e^{-c'J}).
 \label{eq:bulk-Y-readout}
\end{align}
Here $c'>0$ may depend on $c$, and $M_k$ is bounded above and
below by positive constants and is common to all histories,
up to an exponentially small error which is included in the
display.  The coordinate changes also give
\begin{equation}\label{eq:bulk-coordinate-conversion}
 t_k=u_k+O(u_k^2+Z_k^2),\qquad
 z_k=\ell_k^{-1}Z_k+O((|u_k|+Z_k^2)Z_k),\qquad
 \norm{\mathcal R_k}\le C(u_k^2+Z_k^2),
\end{equation}
where $\ell_k$ is the absolute-index fundamental balancing
factor and is bounded above and below positively.
\end{proposition}

\begin{proof}
At step $k$ compare the plane activity with the exact Gaussian
trajectory having parameter $d=u_k$.  The even balanced remainder
has size $O(Z_k^2)$ by Lemma~\ref{lem:bulk-odd}; the odd remainder
has size $O((|u_k|+Z_k^2)Z_k)$.  Undoing the linear balance,
using its parity, gives a difference $O(Z_k^2)$ in the gradient
coordinate and
\begin{equation}\label{eq:bulk-z-conversion}
 z_k=\ell_k^{-1}Z_k+O((|u_k|+Z_k^2)Z_k).
\end{equation}
The Gaussian remainder is $O(u_k^2)$.  These facts prove
\eqref{eq:bulk-coordinate-conversion}.

The stopped torus requires a separate comparison, since its
remainder need not copy the plane.  Run its own exact map from
entry, comparing at intermediate index $r$ with its Gaussian
trajectory at parameter $u_r$.  Gaussian torus trajectories
and their first parameter derivatives are bounded uniformly:
this follows from the same energy completion and contraction
argument as Lemma~\ref{lem:bulk-parity}.  Changing the comparison
parameter from $u_r$ to $u_{r+1}$ costs $O(Z_r^2)$ by
\eqref{eq:bulk-backward-u}.  The remainder forcing has zero
linear singleton derivative at zero.  First the odd torus
contraction gives $O(Z_r)$.  Its even forcing from odd inputs
is then $O(Z_r^2)$.  Comparing even components with the Gaussian
curve and iterating contraction gives $O(Z_r^2)$ in the even
difference.  Returning to the odd equation improves its forcing
to $O((|u_r|+Z_r^2)|Z_r|)$: pure odd terms beyond linear have
at least three odd factors, and every other such term has an
even factor of size $O(|u_r|+Z_r^2)$.  The adjacent-ratio bounds
and \eqref{eq:bulk-backward-u} justify the same geometric
convolutions as in Lemma~\ref{lem:bulk-odd}.  Initial transients
are $O(e^{-c'J})$ on the stated window.

All terminal integrations are analytic on a uniform neighborhood,
with a fixed number $D_0^2$ of blocks and regulator reserve.
In the variance diagnostic a term odd under half-period shift
has zero first derivative at the Gaussian trajectory.  Hence
the difference from that trajectory is $O(Z_k^2)+O(e^{-c'J})$.
It remains to compute the Gaussian value in the coordinate $u$,
rather than approximate a Gaussian trajectory in activity norm.

Hold its initial parameter $d$ fixed.  Let $C_i$ be the reference
tail covariance at the entry side $L^i$.  Direct completion of
the initial quadratic exponential gives, beyond the reference
source prefactor, the log transform
\begin{equation}\label{eq:bulk-Gaussian-readout}
 \frac d2
 \bigl(A^{1/2}F,
 (I-dA^{1/2}C_iA^{1/2})^{-1}A^{1/2}F\bigr).
\end{equation}
On the first mode the operator $A^{1/2}C_iA^{1/2}$ has eigenvalue
$1+O((L^i/N_k)^2)$, by the low-pass polynomial formula.
Together with \eqref{eq:bulk-F-energy}, two source derivatives
therefore give normalized variance
$a/(1-d)+O((L^i/N_k)^c+(s/N_k)^c+N_k^{-c})$.
On a fixed positive-fraction window these errors are
$O(e^{-c'J})$.  Set $d=u_k$ and compare the cell test to the
original sampled cosine in reference energy.  Removing its
noise and cutoff has the same error.  This proves
\eqref{eq:bulk-U-readout}.

For the mean-phase diagnostic the purely even terms integrate
to zero.  Expanding in the odd difference from the Gaussian
trajectory, using \eqref{eq:bulk-z-conversion} and the odd
remainder bound, gives \eqref{eq:bulk-Y-readout}.  The
coefficient at zero activity is the untilted terminal
fundamental multiplier times $\ell_k^{-1}$.  It is real positive
and uniformly bounded above and below.  Finally we verify its
independence from the history.  Its high-mean variance includes
all Gaussian shells below the stop, including those integrated
at entry below $L^{i(J)}$.  These variances telescope to the
same absolute-index sum for every entry choice.  The remaining
tail variance and $\ell_k$ also depend only on $k$.
Observation noise contributes only $O((s/N_k)^2)$ to the
exponent.  Thus a single sequence $M_k$ can be used in
\eqref{eq:bulk-Y-readout}, absorbing the exponentially small
noise correction into its error.  The cutoff removal is again
exponentially smaller still.
\end{proof}

We have proved three facts about every long interior window of a
finite history: its state tends to zero, its odd coordinate is
positive and only subexponentially small, and its two scalar
coordinates can be read from physical observables independent of
entry.  These are the inputs for the trajectory argument that
follows.
